\section{Einsatz im Projekt \emph{Quantum Sphaera}}\label{sec:fallstudie}

Die bisherigen Abschnitte beschreiben QAM ohne Bezug auf ein Anwendungsgebiet. Dieser Abschnitt zeigt, wie das Modell in dem Projekt
eingesetzt wird, für das es entstanden ist. Die Angaben sind Erfahrungsberichte aus einem Einzelfall, keine kontrollierte Studie.

\subsection{Rahmen}

\emph{Quantum Sphaera} ist ein Programm der mathematischen Physik, das von einer Vermutung des Autors ausgeht: dass Materie, Licht und ein
Tachyonensektor ein Objekt sind, das in drei Frames ausgelesen wird~\cite{companion}. Der Autor ist kein Physiker. Sprachmodelle (Claude;
ChatGPT und Codex) wurden von Beginn an eingesetzt und haben den größten Teil des Codes, der Lean-Formalisierung und vieler Texte
geschrieben. Die öffentliche Fassung v5.2.0 vom 17.~November~2026 umfasst 1457 Seiten in der Vollausgabe und 30 Seiten im Physik-Artikel.
Nach einer früheren Zählung bestehen rund \(85\,\%\) des Forschungstexts aus Protokollen, Registern und Prüfrunden. Seit der ersten
geprüften Fassung sind die Runden v4.9.0-r01 bis r86 und v4.10.0-r01 bis r29 hinzugekommen. Im Register der vorab festgeschriebenen Tests
(Vorab-Notizen bis H112) stehen 89 Ergebniszeilen abgeschlossener Tests; einige fassen mehrere Hypothesen zusammen, andere
melden Teilergebnisse. 26 Zeilen lauten auf widerlegt, 17 auf bestätigt und 8 auf ungültig nach dem eigenen Kriterium; die übrigen
38 enthalten gemischte Teilergebnisse, Verträglichkeit, strukturelle Befunde oder Fälle, die keine Entscheidung zuließen.

Die formale Akte deckt davon nur einen kleinen Teil ab. Im Repository liegen zehn exportierte \texttt{.qam}-Akten mit zusammen
81 Einträgen, vor allem Beispiele und zwei Runden. Der weitaus größere Teil des Projekts benutzt QAM als \emph{Disziplin}: Statussprache,
Guard-Listen und Register in Notizen und Runden, Belege der Maschine für Rechnungen. Das ist eine wichtige Einschränkung jeder Aussage
über die \glqq Skalierung\grqq{} auf rund 1450 Seiten: Skaliert hat die Arbeitsweise, nicht die mechanisch geprüfte Akte.

\subsection{Wie QAM im Projekt benutzt wird}

\begin{itemize}[nosep]
\item \emph{Append-only-Runden.} Neue Ergebnisse kommen als nummerierte Runden in eigene Schichtdateien. Frühere Runden werden nie
umgeschrieben; eine Korrektur ist eine neue Runde, die nennt, was sie korrigiert. Das ist (R1) auf der Ebene des Forschungstexts.
\item \emph{Vorab-Notizen.} Jeder Test mit Entscheidungskriterium wird vor der Rechnung mit Hypothese, Daten und Kriterium festgeschrieben
und committet; der Zeitstempel des Commits fixiert ihn~\cite{nosek}. Ergebnisse, die vor der Rechnung analytisch feststanden, werden als
\glqq nicht blind\grqq{} geführt. Eine Vorab-Notiz hat in QAM-Begriffen die Rolle einer These mit festem Wortlaut.
\item \emph{Status in Prüfnotizen.} Mathematische Selbstprüfungen schließen mit einem Urteil in der Statussprache, etwa \glqq proved as
written\grqq{} für eine bedingte Implikation und \glqq conditional on a named input\grqq{} für ihre Anwendung, und nennen die offenen Guards.
\item \emph{Gegenterm-Register.} Fehlende Begründungen werden als nummerierte Platzhalter mit eigenen Guards geführt. Vorhersagen, die
auf einem Platzhalter ruhen, erben dessen Guards. So sind drei Vorhersagen (H15, H46, H96) mit dem Platzhalter CT-004 verknüpft und tragen
dessen vier offene Bedingungen; ihre Prüfkriterien ändern sich dadurch nicht, aber ihr Status bleibt der einer Lesart.
\item \emph{Rechnungen als Aufträge.} Suchen und exakte Rechnungen laufen als MML-Aufträge; Ergebnisse werden per Belegshash zitiert.
Gewichtige Ergebnisse werden mit unabhängig geschriebenem Code nachgerechnet, mit Computeralgebra (Wolfram Language) für symbolische und
exakte Rechnungen und mit Qiskit für explizit definierte Schaltungen.
\end{itemize}

\subsection{Konstruiertes Beispiel: ein Projektor und die Koide-Relation}

Das folgende \glqq Vorher\grqq{} ist \emph{konstruiert}. Es gibt einen typischen unzulässigen Schluss wieder, ist aber nicht das Protokoll
eines tatsächlichen Gesprächs. Das \glqq Nachher\grqq{} beruht auf den dokumentierten Prüfnotizen des Projekts.

\paragraph{Hintergrund.} Die Koide-Relation~\cite{koide} verknüpft die Massen der geladenen Leptonen,
\[
K=\frac{m_e+m_\mu+m_\tau}{\bigl(\sqrt{m_e}+\sqrt{m_\mu}+\sqrt{m_\tau}\bigr)^2}\approx\frac23 .
\]
Im Projekt tritt ein Kompensator \(C=(I+J)/2\) auf, wobei \(J\) eine orthogonale Involution ist. Wegen \(J^2=I\) gilt \(C^2=C\), und mit
\(J^*=J\) ist \(C\) ein orthogonaler Projektor.

\paragraph{Vorher (konstruiert).} Ein frei schreibendes Modell könnte formulieren:
\begin{quote}\itshape
\glqq Der Kompensator \(C=(I+J)/2\) ist ein Projektor. Wendet man ihn auf den symmetrischen Wurzelvektor an, ergibt sich \(K=2/3\). Die
Koide-Relation folgt also aus der Kompensatorstruktur.\grqq
\end{quote}
Jeder Satz klingt plausibel, und der erste ist richtig. Der Schluss enthält aber mindestens zwei stille Übergänge.

\paragraph{Was tatsächlich gilt.} Auf einem reellen Wurzelraum \(\mathbb R^3\) sei \(e=(1,1,1)/\sqrt3\), \(J_m\) eine orthogonale Involution,
\(C_m=(I+J_m)/2\) und \(x=\alpha\,C_m e\) mit \(\alpha>0\) ein nichtverschwindender, komponentenweise nichtnegativer Vektor. Mit
\(h=e\cdot J_m e\) gilt \(\lVert C_m e\rVert^2=(1+h)/2\), die Komponentensumme ist \(\alpha\sqrt3\,(1+h)/2\), und daraus
\[
K=\frac{\lVert x\rVert^2}{\bigl(\sum_i x_i\bigr)^2}=\frac{2}{3(1+h)},\qquad h>-1 .
\]
\(K=2/3\) gilt also genau dann, wenn \(h=0\). Für \(J_m=I\) ist \(h=1\) und \(K=1/3\). Ein zweites Ergebnis überträgt die Struktur aus einem
Hilbertraum \(H\) mit Involution \(J\) über eine Koisometrie \(B\colon H\to\mathbb R^3\) mit \(J\)-invariantem Kern: Dann ist \(J_m=BJB^*\) eine
orthogonale Involution und \(BC=C_mB\). Auch hier liefert \(J=I\), \(B=I\) den Wert \(K=1/3\); der Abstieg allein erzwingt Koide nicht. Beide
Aussagen wurden in den Prüfnotizen als \glqq proved as written\grqq{} bestätigt.

Der freie Schluss setzt damit stillschweigend voraus:
\begin{enumerate}[nosep]
\item die \emph{Orthogonalitätsbedingung} \(h=0\), die aus der Projektoreigenschaft nicht folgt;
\item die \emph{Zuordnung}: dass es eine native Abbildung \(B\) und einen unabhängig gelieferten Startvektor gibt, dass die Komponenten von
\(x\) die Wurzeln physikalischer Massen sind und positiv.
\end{enumerate}
Die erste ist eine mathematische Zusatzannahme, die zweite eine physikalische Lesart.

\paragraph{Nachher.} In QAM wird dieselbe Überlegung als Akte geführt. Der folgende Block steht im Quelltext dieses Artikels und wird beim
Setzen geprüft:
\begin{lstlisting}
\qamframe{H} \qamframe{W}                 % Hilbertraum, Wurzelraum
\qammap[B]{B}{H}{W}
\qamclaim[D]{descent}{H}{proved}           % J_m = B J B*, BC = C_m B
\qamconvert[D']{descentW}{proved}{B}{descent}
\qamclaim[K_h]{koidehc}{W}{proved}         % K = 2/[3(1+h)]
\qaminput[\nu]{native}{W}{reading}         % B nativ, Startvektor, Wurzelgesetz, Positivität
\qaminput[o]{orth}{W}{conditional}         % h = 0
\qamderive[K]{koide}{reading}{descentW,koidehc,native,orth}
\end{lstlisting}
\qamframe{H}\qamframe{W}
\qammap[B]{B}{H}{W}
\qamclaim[D]{descent}{H}{proved}
\qamconvert[D']{descentW}{proved}{B}{descent}
\qamclaim[K_h]{koidehc}{W}{proved}
\qaminput[\nu]{native}{W}{reading}
\qaminput[o]{orth}{W}{conditional}
\qamderive[K]{koide}{reading}{descentW,koidehc,native,orth}
Das Paket setzt die Koide-Aussage als
\begin{center}\qam{koide}\end{center}
und \verb|\qamshow{koide}| ergibt \qamshow{koide}. Hätte der Quelltext stattdessen
\verb|\qamderive{koide}{proved}{...}| gelautet, also den freien Schluss nachgebildet, wäre die Übersetzung mit einer Beförderungsmeldung abgebrochen:
\begin{lstlisting}
! Package qam Error: Promotion: `koide' is declared `proved', but its weakest
(qam)                premise `native' is only `reading'.
\end{lstlisting}
Die beiden bewiesenen Teile bleiben bewiesen und sind wiederverwendbar; die Koide-Aussage trägt ihre beiden offenen Bedingungen als
Eingaben, und jeder spätere Schluss aus ihr erbt sie.

Das Beispiel zeigt auch eine Grenze. Dass \(B\) eine Abbildung zwischen den Frames ist, kann in der Akte nur deklariert werden; Abbildungen
haben keinen Status. Die tatsächliche Prüfnotiz des Projekts hält deshalb fest, dass der Export weiterhin \texttt{frame-bridge-required}
meldet, weil es noch keinen maschinell geprüften Typ für solche Abbildungen gibt, und dass diese Meldung nicht durch Umbenennen der Frames
unterdrückt wird. Die offenen Guards der Notiz lauten: native Abbildung und Startvektor mit Wurzelgesetz, unabhängig begründete
Orthogonalität, maschinenprüfbare typisierte Frame-Abbildung.

\subsection{Protokollierter Fall: eine Eingabe als Ergebnis}

Ein zentraler Parameter des Projekts, \(\beta\approx0{,}549\), war aus zwei Literaturwerten eingesetzt worden und wurde später als
hergeleitet zitiert. Eine Formel, die ihn auf \(0{,}018\,\%\) traf, hatte einen nachträglich gewählten Exponenten. In QAM ist \(\beta\) die
Auslesung eines Zeugen, \({}^{M}\beta=R_\beta(u)\) mit \(u=\langle6,(\gamma_1,\gamma_2)\rangle\); jede Aussage, die auf \(\beta\) ruht, trägt
\(\gamma_1,\gamma_2\) in ihrer Eingabemenge, und eine Unabhängigkeitserklärung für sie hält den Compiler an. Die heutige Vorhersage zur
Materiedichte, \(\Omega_M=\beta^2\), wird deshalb ausdrücklich als zusätzliche Gleichsetzung und mit dem Status einer Lesart geführt. Eine
spätere Vorhersage (H112) trennt die These \glqq beide Messgrößen zeigen dieselbe Zahl\grqq{} vom eingesetzten Wert.

\subsection{Was gefangen wurde und was nicht}

\begin{center}
\small
\renewcommand{\arraystretch}{1.25}
\begin{longtable}{@{}>{\raggedright\arraybackslash}p{0.19\textwidth}>{\raggedright\arraybackslash}p{0.45\textwidth}>{\raggedright\arraybackslash}p{0.28\textwidth}@{}}
\toprule
Fehlerklasse & Beispiel aus dem Projekt & aufgefangen durch\\
\midrule
\endhead
Frame-Vermischung & Ein \glqq struktureller Wert\grqq{} setzte eine Auslesung eines Frames mit einem in einem anderen Frame gemessenen Versatz gleich & Regel (R4): keine deklarierte Abbildung\\
Eingabe als Ergebnis & \(\beta\) als hergeleitet zitiert (oben) & Zeugentypisierung, Regel (R3)\\
Beförderung & Bedingte oder lokale Aussagen als Sätze geführt, etwa eine Kette zur kosmologischen Konstante & Status, (R2), Prüfrunden\\
externe Abhängigkeit & Beweise, die eine separate Arbeit mit behauptetem Beweis der Riemannschen Vermutung zitierten & \emph{nicht} durch QAM, sondern durch eine Prüfrunde\\
numerisches Artefakt & Ein Hochpräzisionsergebnis aus 100-stelligen Zahlen, die durch eine 15-stellige Textumwandlung liefen & Nachrechnung an einem generischen Punkt\\
Literaturlesart & Ein Zählwert als \glqq korrigiert\grqq{} gelesen, den spätere Literatur für weniger gut begründet hält & unabhängige Literaturprüfung\\
\bottomrule
\end{longtable}
\end{center}

Die drei unteren Zeilen sind die wichtigere Hälfte der Tabelle. QAM hat keinen dieser Fehler gefunden und konnte es nicht: Eine falsch
zitierte Quelle, ein Rundungsfehler im Werkzeug und eine fragwürdige Lesart der Literatur sind inhaltliche Fehler. Die Struktur der Akte
zeigt nur, \emph{worauf} eine Aussage ruht; ob das Fundament trägt, muss anders geprüft werden.

\subsection{Erfahrungen}

Drei Beobachtungen aus dem Projekt verallgemeinern wir vorsichtig. Erstens kostet die Disziplin Umfang: Der größte Teil des Texts ist
Protokoll. Zweitens verschiebt sie die Arbeit der Sprachmodelle: Statt Behauptungen zu formulieren, formulieren sie Aufträge, Prüfnotizen
und Korrekturrunden, und ein großer Teil ihrer Arbeit besteht darin, offene Bedingungen genau zu benennen. Drittens wirkt die Akte
vor allem gegen schleichende Verschiebungen über viele Sitzungen: Ein Modell, das eine frühere Bedingung vergessen hat, findet sie im
Status des Eintrags wieder, den es wiederverwenden will. Diese Beobachtungen sind qualitativ; den Zeitaufwand der Buchführung haben wir
nicht gemessen.
